\subsection{Large MuDst sample: method 1 (emulation)}
\label{sec:big}

\paragraph{Sample and production.} Restored from HPSS (\fpath{scripts/mudst_restore100.sh}, list
\fpath{filelists/mudst/restore_100.txt}): 98 MuDst files, 7 per run for the 14 runs of the picoDst sample without 20088015.
Together with the four files of Sec.~\ref{sec:mudst} the sample has 102 files (\fpath{filelists/mudst/tags_all.txt}) and
\mubigEvAll{} events, of which \mubigEvMtwo{} pass the selection of Sec.~\ref{sec:m2} and are used for $m^2$. One HTCondor job
per file:
\begin{enumerate}
\item calibration pass: \fpath{macros/probeMuBTof.C} pass 1 on the first 4000 events of each file
  (\fpath{scripts/probe1.sub}); the outputs of a run are merged into the per-cell offsets \fpath{output/mudst/off_<run>.root}
  (\fpath{scripts/off_big.sh});
\item method 1, \fpath{StTofMEMuMaker} v4 (\fpath{scripts/mubig.sub}, \fpath{scripts/mu_file.sh}): production emulation with
  partner-track claims, as in Sec.~\ref{sec:muemul}. Mixing buffers are per run and file;
\item method 2, \fpath{StTofMEOverlayMaker} v3 with partner tracks (\fpath{scripts/ovbig.sub}); results in
  Sec.~\ref{sec:bigm2}.
\end{enumerate}
Both makers (v4 and v3) additionally store the charge-signed $q\times m^2$ and the $m^2$ of negative tracks for SE and ME; the
positive-track histograms are the difference. The 102 method-1 outputs are merged into \fpath{data/mu_big.root}. Scripts:
\fpath{scripts/make_mu.C} (tag \fcmd{big}), \fpath{scripts/make_qm2.C}, \fpath{scripts/make_fit.C},
\fpath{scripts/make_big_tables.py}; logs \fpath{logs/mubig_numbers.txt}, \fpath{logs/qm2_mubig.txt},
\fpath{logs/fit_mubig*.txt}. Statistical uncertainties only.

\caveat{All numbers of this section are from method 1, which uses our own hit times (RMS \mubigCalRMS\,ns against production;
Sec.~\ref{sec:muemul}). Method 2 and the comparison follow in Sec.~\ref{sec:bigm2}.}

\subsubsection{Random matches}
The emulation picks the production hit for \mubigChoiceSame\% of the production matches, as in the four-file test. Random
matches are \mubigFrProdCi\% of the TOF matches in 0--5\% and \mubigFrProdAll\% in 0--80\% ($0.2<p<3$\,GeV/$c$), consistent with the
four-file values (\muFrProdCi\% and \muFrProdAll\%). Between the p and d peaks SE/ME is \mubigHiProdCi{} (0--5\%) and
\mubigHiProdAll{} (0--80\%). Per momentum bin (Table~\ref{tab:bigbins}) it is \bigHiMidMin--\bigHiMidMax{} at 0.6--2\,GeV/$c$
(0--5\%). The excess that remains is at $0.2<p<0.4$\,GeV/$c$
(\bigHiLow) and above 2\,GeV/$c$ (\bigHiTopMin--\bigHiTopMax). The largest bin-by-bin ME/SE values are in nearly empty bins just
below the pion peak; they are the timing artefact of method 1 (red box in Sec.~\ref{sec:muemul}).

\begin{figure}[H]
\centering
\figpanel{0.9}{mubig_m2_cent}{$m^2$ of QA tracks, 0--5\%, nine momentum bins, 102 MuDst files (log scale). Black points: same event
(production $\beta$). Red filled: mixed event, method 1 with partner-track claims, normalized per QA track. Random matches \mubigFrProdCi\%; SE/ME between p and d \mubigHiProdCi.}
\caption{Large sample, method 1, SE vs ME, 0--5\% (\fpath{scripts/make_mu.C}).}
\end{figure}

\input{tables/big_m1_bins}

\paragraph{Other centralities.} Figures~\ref{fig:mubig_m2_mid} and \ref{fig:mubig_m2_per} show the same comparison for mid-central
(30--40\%) and peripheral (60--80\%) events, and Table~\ref{tab:bigcent} gives the numbers per momentum bin. Random matches fall
with centrality: \mubigFrProdCi\% (0--5\%), \bigFrMid\% (30--40\%) and \bigFrPer\% (60--80\%) of the TOF matches at $0.2<p<3$\,GeV/$c$.
The mixed event describes the region between the p and d peaks best in central events; SE/ME there is \mubigHiProdCi{} (0--5\%),
\bigHiMid{} (30--40\%) and \bigHiPer{} (60--80\%). In the lowest momentum bin it is \bigHiLow, \bigHiLowMid{} and \bigHiLowPer: the
SE continuum above the ME at $0.2<p<0.4$\,GeV/$c$ is seen in all three centralities and is relatively largest in peripheral events,
where the random background is smallest (read from the figures: in 60--80\% the SE continuum at 0.2--0.6\,GeV/$c$ lies several times
above the ME, cf.\ the HI column of Table~\ref{tab:bigcent}). In 60--80\% above 1.2\,GeV/$c$ the ME between p and d is close to zero and the HI ratio is set by a
few SE counts.

\caveat{A source of SE entries above the ME that scales with the number of tracks rather than with the hit density (e.g.\ late hits
correlated with the track itself) would give exactly this centrality pattern. It is studied in Sec.~\ref{sec:lowp}.}

\begin{figure}[H]
\centering
\figpanel{0.9}{mubig_m2_mid}{$m^2$ of QA tracks, 30--40\%, nine momentum bins, 102 MuDst files (log scale). Black points: same event;
red filled: mixed event, method 1 with partner-track claims. Random matches \bigFrMid\%; SE/ME between p and d \bigHiMid.}
\caption{Large sample, method 1, SE vs ME, 30--40\% (\fpath{scripts/make_mu.C}).}
\end{figure}

\begin{figure}[H]
\centering
\figpanel{0.9}{mubig_m2_per}{$m^2$ of QA tracks, 60--80\%, nine momentum bins, 102 MuDst files (log scale). Encodings as in
Fig.~\ref{fig:mubig_m2_mid}. Random matches \bigFrPer\%; SE/ME between p and d \bigHiPer.}
\caption{Large sample, method 1, SE vs ME, 60--80\% (\fpath{scripts/make_mu.C}).}
\end{figure}

\input{tables/big_m1_cent}

\subsubsection{Charge-separated $q\times m^2$}
Fig.~\ref{fig:qm2_mubig_cent} shows SE, ME and SE$-$ME versus $q\times m^2$, which separates $\pi^\pm$, K$^\pm$ and p/$\bar{\rm p}$;
the deuteron appears on the positive side only (read by eye). The random fraction is similar for both charges at low momentum and
larger for negative tracks at high momentum ($f^-$, $f^+$ in Table~\ref{tab:bigbins}). The random hits do not depend on the
track charge; the fraction is larger for negatives because there are fewer real heavy negative particles to normalize to.

\begin{figure}[H]
\centering
\figpanel{0.9}{qm2_mubig_cent}{$q\times m^2$ of QA tracks, 0--5\%, 102 MuDst files, 0.02\,GeV$^2$/$c^4$ bins (log scale). Black line:
same event; red filled: mixed event (method 1); blue points: SE$-$ME.}
\caption{Large sample, method 1, charge-signed $m^2$, 0--5\% (\fpath{scripts/make_qm2.C}).}
\end{figure}

\subsubsection{Signal extraction: SE$-$ME and line-shape fits}
\label{sec:bigfit}
In each centrality group and momentum bin, the SE histogram (0.01\,GeV$^2$/$c^4$ bins) is fitted in $-0.6<m^2<4.8$\,GeV$^2$/$c^4$
with a Poisson likelihood. The model is signal + constant + the normalized ME, with the ME held fixed; this is the SE$-$ME
subtraction with correct statistics. The function is integrated over each bin. The signal has four peaks ($\pi$, K, p, d), and
four line shapes are compared, each with free yield, mean and width per species:
\begin{itemize}
\item G: Gaussian;
\item G2: double Gaussian (common mean, core and wide width, core fraction);
\item T: Student-t, $A/\sigma\,t_\nu((m^2-\mu)/\sigma)$, free $\nu$ per species;
\item T2: split Student-t, separate widths below and above the mean, free $\nu$ per species.
\end{itemize}
The deuteron width is bounded so that it cannot absorb a continuum. The model with the smallest Pearson $\chi^2$/ndf of SE
against the model is called best.

\choice{The Student-t family describes the peaks better than the Gaussian family. The split Student-t (T2) is the best model in
\bigFitTtwoBestCi{} of \bigFitNbinsCi{} momentum bins in 0--5\% and \bigFitTtwoBestAll{} of \bigFitNbinsAll{} in 0--80\% (all
tracks). Mean $\chi^2$/ndf in 0--5\%: G \bigFitChiGCi, G2 \bigFitChiGtwoCi, T \bigFitChiTCi, T2 \bigFitChiTtwoCi. The asymmetric
shape is needed for the high-$m^2$ tail of the proton peak. T2 is used for the yields.}

\begin{figure}[H]
\centering
\figpanel{0.9}{fit_mubig_cent}{SE$-$ME $m^2$ (points), 0--5\%, all tracks, 102 MuDst files, method 1 (log scale), with the signal part
of the best Gaussian-family model (blue dashed) and the best Student-t-family model (red). The $\chi^2$/ndf of both are given
in each panel.}
\caption{Large sample, method 1: SE$-$ME and line-shape fits, 0--5\% (\fpath{scripts/make_fit.C}).}
\end{figure}

\input{tables/big_m1_fit}

The fitted yields add up to SE$-$ME within \bigSumYDev\% in every bin. With T2, $\chi^2$/ndf is \bigTtwoChiTopMin--\bigTtwoChiTopMax{}
above 2\,GeV/$c$, \bigTtwoChiHiMin--\bigTtwoChiHiMax{} at 1--2\,GeV/$c$, \bigTtwoChiMidMin--\bigTtwoChiMidMax{} at 0.4--1\,GeV/$c$ and
\bigTtwoChiLowP{} at 0.2--0.4\,GeV/$c$ (0--5\%).

\caveat{No model is statistically acceptable below 2\,GeV/$c$ with \mubigEvMtwo{} events. Two limitations are seen in the figures (read by
eye): (i) at $0.2<p<0.4$\,GeV/$c$ a continuum of about 100 counts per bin remains between the peaks after the subtraction, the same
SE excess as in Table~\ref{tab:bigbins}; it is under-subtracted background, not a line-shape problem; (ii) each momentum bin mixes
tracks with different $m^2$ resolution (momentum, path length, centrality), so a single line shape per species cannot be exact at this
precision. Finer momentum bins or a resolution-scaled variable ($1/\beta$ residuals) would address (ii).}

\paragraph{Charge-separated fits.} The same fits were done separately for negative and positive tracks
(Figs.~\ref{fig:fit_mubigneg_cent} and \ref{fig:fit_mubigpos_cent}). The Student-t family is better in \bigFitNegTfamCi{} (negative)
and \bigFitPosTfamCi{} (positive) of 9 bins in 0--5\%. The sum of the negative and positive yields reproduces the all-track yield
within \bigClosPiDev\% ($\pi$), \bigClosKDev\% (K) and \bigClosPDev\% (p) at 0.4--2.5\,GeV/$c$. In that range the ratios are
$\pi^-/\pi^+ = $ \bigRatPiMin--\bigRatPiMax, K$^-$/K$^+ = $ \bigRatKMin--\bigRatKMax{} and $\bar{\rm p}$/p $=$ \bigRatPMin--\bigRatPMax{}
(0--5\%, Table~\ref{tab:bigratio}); these are raw yield ratios without efficiency corrections. At 0.2--0.4\,GeV/$c$ the
$\bar{\rm p}$/p ratio is distorted by the continuum, and at 2.5--3\,GeV/$c$ the $\pi$ and K peaks are no longer separated.

\input{tables/big_m1_ratio}

\begin{figure}[H]
\centering
\figpanel{0.82}{fit_mubigneg_cent}{Negative tracks: SE$-$ME $m^2$ and fits, 0--5\%; encodings as in Fig.~\ref{fig:fit_mubig_cent}.}
\caption{Large sample, method 1: SE$-$ME and line-shape fits, negative tracks, 0--5\%.}
\end{figure}

\begin{figure}[H]
\centering
\figpanel{0.82}{fit_mubigpos_cent}{Positive tracks: SE$-$ME $m^2$ and fits, 0--5\%; encodings as in Fig.~\ref{fig:fit_mubig_cent}.}
\caption{Large sample, method 1: SE$-$ME and line-shape fits, positive tracks, 0--5\%.}
\end{figure}
